\documentclass[10pt,a4paper]{amsart}
\usepackage[margin=19mm]{geometry}
\usepackage{amsmath,amssymb,mathtools}
\usepackage[scr=rsfs,cal=euler]{mathalfa}
\usepackage{xfrac}
\usepackage[unicode,hidelinks,bookmarks=false]{hyperref}
\newcommand{\ie}{i.e.\ }
\newcommand{\cf}{cf.\ }
\newcommand{\resp}{respectively\ }
\newcommand{\Subsection}{\subsection}
\providecommand{\nobreakdash}{\nobreak\hbox{-}}
\setlength{\emergencystretch}{2em}
\multlinegap0pt
\usepackage{bm}
\usepackage{bbm}
\usepackage{dsfont}
\usepackage{xspace}
\usepackage{cancel}
\usepackage{mathtools}
\usepackage{cleveref}
\usepackage{amsthm}
\makeatletter
\@ifundefined{theorem}{\newtheorem{theorem}{Theorem}}{}
\@ifundefined{claim}{\newtheorem{claim}[theorem]{Claim}}{}
\@ifundefined{lemma}{\newtheorem{lemma}[theorem]{Lemma}}{}
\makeatother
\newcommand{\E}[1]{\mathbb{E}\left[{#1}\right]}
\newcommand{\mc}{\preceq_{\textnormal{mc}}}
\newcommand{\nug}{\nu^{\textnormal{GD}}}
\newcommand{\nus}{\nu^{\textnormal{sim}}}
\newcommand{\psim}{P_{\textnormal{s-GD}}}
\newcommand{\sd}{\preceq_{\textnormal{sd}}}
\title[Rapid Mixing of Glauber Dynamics for Monotone Systems via En]{Rapid Mixing of Glauber Dynamics for Monotone Systems via Entropic Independence}
\author{Weiming Feng, Minji Yang}
\date{Source: arXiv:2507.11031v1 (CC BY 4.0)}
\begin{document}
\maketitle

\noindent\emph{Source and licence.} Rapid Mixing of Glauber Dynamics for Monotone Systems via Entropic Independence. Version arXiv:2507.11031v1 (CC BY 4.0), distributed under the Creative Commons Attribution 4.0 International licence, as recorded in the arXiv licence metadata for this exact version and shown on the abstract page https://arxiv.org/abs/2507.11031v1. Full rights notice: licenses/arxiv-2507.11031-CC-BY-4.0.txt.\par\smallskip

\noindent\textit{Editorial scope.} Counted unit: the Lemma whose source label is \texttt{lem:phase-comparison-single-vertex} in the article (source0.tex lines 1178--1180, source label \texttt{lem:phase-comparison-single-vertex}), delivered together with the article's own complete original proof of it (source0.tex lines 1181--1274, one \texttt{proof} environment of 94 lines). No mathematics is written, repaired or re-derived: statement and proof are the article's own text line for line. Required context retained uncounted: none. Every definition, display and label of those lines is retained unchanged. Omitted from this package and not redistributed: 1--1177, 1275--2211. Nothing inside a retained excerpt is omitted or altered: every retained body is delivered exactly as the article's lines are written, so no omission receipt of any kind accompanies this package. Citations: the retained excerpts carry no citation command at all, so no bibliography entry of the article is transcribed and no reference is redistributed. Reference adaptation: none was needed and none was made; every reference inside the delivered unit resolves inside it, so no unresolved reference or citation is produced. Printed slips kept literal: no printed word, symbol, hypothesis or number of the counted statement or of its proof was altered. Layout and class: the article's own class is \texttt{article} with packages including amsmath, amssymb, amsthm, thmtools, hyperref, cleveref, fullpage, microtype, caption, bbm, color, algorithm2e, bm, natbib; the wrapper is the lane's pinned \texttt{amsart} editorial wrapper at 10pt on A4, which restates the theorem-like environments the retained text needs and adds the article's own macro definitions that the retained text needs (\texttt{\textbackslash E}, \texttt{\textbackslash mc}, \texttt{\textbackslash nug}, \texttt{\textbackslash nus}, \texttt{\textbackslash psim}, \texttt{\textbackslash sd}), with extra wrapper packages (bm, bbm, dsfont, xspace, cancel, mathtools, cleveref). The delivered numbering is the unit's own. Assets: no retained span contains a figure environment, a graphics-inclusion command, a PDF-page inclusion, a table or a TikZ picture; no image file is copied into this package, the package declares no asset dependency and no image file is redistributed with it. Licence: version arXiv:2507.11031v1 of this article is distributed under the Creative Commons Attribution 4.0 International licence, as recorded in the arXiv licence metadata for this exact version and shown on the abstract page https://arxiv.org/abs/2507.11031v1. A full rights notice is delivered at licenses/arxiv-2507.11031-CC-BY-4.0.txt. Characters: every retained excerpt is pure ASCII, so the delivered engine, XeLaTeX with its default Latin Modern text font, reports no missing character.
\begin{lemma}\label{lem:phase-comparison-single-vertex}
For any vertex $i \in V$, $P^i_{\pi\textnormal{-GD}} \mc \psim^i$.
\end{lemma}

\begin{proof}
%Let $P_\text{alg}^i$ denote the transition in $P_\text{alg}$ conditioned on the vertex $i \in V$ is being selected in the first step. Let $P^i_{\pi\textnormal{-GD}}$ denote the transition in $P_{\pi\textnormal{-GD}}$ conditioned on the vertex $i \in V$ is being selected by the Glauber dynamics. We view both $P_\text{alg}^i$ and $P^i_{\pi\textnormal{-GD}}$ as the transition matrices in space $\{0,1,\star\}^V$. We have 
Recall that $\Omega(\pi)$ is the support of $\pi$.
Fix a vertex $i \in V$. 
By the definition of $\mc$, we need to show that for any distribution $\nu$ over $\Omega(\pi)$ such that the function $\frac{\nu}{\pi}:\Omega(\pi) \to \mathbb{R}_{\geq 0}$ is increasing, $\nu P^i_{\pi\textnormal{-GD}} \preceq_{\text{sd}} \nu \psim^i$. Define two distributions
\begin{align*}
\nug = \nu P_{\pi\textnormal{-GD}}, \quad \nus = \nu \psim^i.
\end{align*}
To show $\nug \sd \nus$, it suffices to show that for any increasing function $f:\{0,1,\star\}^V\to \mathbb{R}_{\geq 0}$\footnote{To prove the stochastic dominance, we only need to consider the increasing functions $f$ that are supported on $\Omega(\pi)$. Here, we consider a more general class of functions over $\{0,1,\star\}^V$, which is sufficient for the proof. The benefit is that in the rest of the proof (such as~\eqref{eq:exp-nu-i-1}), $f(\sigma)$ is well-defined for any $\sigma \in \{0,1,\star\}^V$.},
\begin{align}\label{eq:ineq-i-each}
    \E[\nug]{f}\leq \E[\nus]{f}.
\end{align}






For $x\in \{0,1,\star\}$ and a configuration $\sigma\in \{0,1,\star\}^V$, define $\sigma^x$ to be a configuration by changing the value of $\sigma_i$ to $x$, and keeping other values the same as $\sigma$. Formally,
\begin{align*}
    \begin{cases}
&\sigma^x_{i} = x,\\
&\sigma^x_j = \sigma_j,\quad \forall j\in V\setminus \{i\}.   
\end{cases}
\end{align*}





By the definition of $P^i_{s\textnormal{-GD}}$,
we have the following relation between $\nus$ and $\nu$. For any $\sigma\in \{0,1,\star\}^V$ such that $\sigma_{V \setminus \{i\}}$ is a feasible partial configuration on $V \setminus \{i\}$, it holds that
\begin{align}\label{eq:recur-alg}
\begin{split}
&\nus(\sigma^{\star}) = \nu(\sigma^{\star}),\\
&\nus(\sigma^1)=(\nu(\sigma^0)+\nu(\sigma^1))\frac{\theta \mu(\textsf{contr}(\sigma^1))}{\mu(\textsf{contr}(\sigma^0))+\theta \mu(\textsf{contr}(\sigma^1))},\\
&\nus(\sigma^0)=(\nu(\sigma^0)+\nu(\sigma^1))\frac{\mu(\textsf{contr}(\sigma^0))}{\mu(\textsf{contr}(\sigma^0))+\theta \mu(\textsf{contr}(\sigma^1))}.
\end{split}
\end{align}
We explain the above three equations in the following way. Let $X \sim \nu$ and $X'$ be the configuration obtained by applying $P^i_{s\textnormal{-GD}}$ to $X$.
The first equation is because $X'_i = \star$ if and only if $X_i = \star$.
For the second and the third equations, if $X' = \sigma^0$ or $\sigma^1$, then $X_i = 0$ or $1$ and $X'_i$ is resampled from $(\theta * \mu)_i^\tau$ for $\tau = \textsf{contr}(\sigma)_{V \setminus \{i\}}$. Note that in the contracting operation, the value $\star$ is mapped to $1$. We can calculate
\begin{align}\label{eq:recur-alg-2}
    (\theta * \mu)_i^{\tau}(1) = \frac{(\theta * \mu)(\textsf{contr}(\sigma^1))}{(\theta * \mu)(\textsf{contr}(\sigma^0))+(\theta * \mu)(\textsf{contr}(\sigma^1))}.
\end{align}
Let $\Vert \sigma \Vert_1$ be the number of $1$s in $\sigma$. Note that $\Vert \textsf{contr}(\sigma^1) \Vert_1 = \Vert \textsf{contr}(\sigma^0) \Vert_1 + 1$. Hence, the second and the third equations in~\eqref{eq:recur-alg} hold by~\eqref{eq:recur-alg-2}.

Consider the expectation $\E[\nus]{f}$. It can be written as
\begin{align}\label{eq:exp-nu-i-1}
\E[\nus]{f}=\sum_{\sigma\in \{0,1,\star\}^V} \nus(\sigma)f(\sigma)=\frac{1}{3} \sum_{\sigma\in \{0,1,\star\}^V}\sum_{x\in\{0,1,\star\} } \nus(\sigma^x)f(\sigma^x). 
\end{align}
The last equation holds because we enumerate all full configurations $\sigma \in \{0,1,\star\}^V$. For $\sigma = \sigma^{x}$ where $x\in \{0,1,\star\}$, the second sum over $x$ is the same. Hence, any $\sigma \in \{0,1,\star\}^V$ is counted $3$ times in the last summation.
Note that for the last sum over $\sigma$, we only need to consider $\sigma$'s that $\sigma_{V \setminus \{i\}}$ is a feasible partial configuration on $V \setminus \{i\}$. Otherwise, $\pi(\sigma^x) = 0$ for all $x\in \{0,1,\star\}$. Since Markov chain $P^i_{s\textnormal{-GD}}$ always stays in feasible space, $\nus(\sigma^x) = 0$ for all $x$.
Fix an arbitrary one of such configurations $\sigma \in \{0,1,\star\}^V$. Consider the term $\sum_{x\in \{0,1,\star\}} \nus(\sigma^x)f(\sigma^x)$ in~\eqref{eq:exp-nu-i-1}.
Using the relation in~\eqref{eq:recur-alg}, we have
\begin{align}\label{eq:exp-nu-i}
&\sum_{x\in \{0,1,\star\}} \nus(\sigma^x)f(\sigma^x)= \nu(\sigma^\star) f(\sigma^\star)+(\nu(\sigma^0)+\nu(\sigma^1))\frac{\theta \mu(\textsf{contr}(\sigma^1))f(\sigma^1)}{\mu(\textsf{contr}(\sigma^0))+\theta \mu(\textsf{contr}(\sigma^1))}\notag\\
&\qquad\qquad\qquad\qquad\qquad+(\nu(\sigma^0)+\nu(\sigma^1))\frac{ \mu(\textsf{contr}(\sigma^0))f(\sigma^0)}{\mu(\textsf{contr}(\sigma^0))+\theta \mu(\textsf{contr}(\sigma^1))} \notag\\
&=\nu(\sigma^\star) f(\sigma^\star)+(\nu(\sigma^0)+\nu(\sigma^1))\frac{ \mu(\textsf{contr}(\sigma^0))f(\sigma^0)+\theta \mu(\textsf{contr}(\sigma^1))f(\sigma^1)}{\mu(\textsf{contr}(\sigma^0))+\theta \mu(\textsf{contr}(\sigma^1))}.  
\end{align}

Define function $g:\Omega(\pi)\to \mathbb{R}$ by: for any $\sigma\in\Omega(\pi)$,
\begin{align}\label{def:g-i}
\begin{split}
&g(\sigma^\star)=f(\sigma^\star);\\
&g(\sigma^0)=g(\sigma^1)=\frac{ \mu(\textsf{contr}(\sigma^0))f(\sigma^0)+\theta \mu(\textsf{contr}(\sigma^1))f(\sigma^1)}{\mu(\textsf{contr}(\sigma^0))+\theta \mu(\textsf{contr}(\sigma^1))}. 
\end{split}
\end{align}
The definition of $g$ depends only on $f$. Then
\begin{align}\label{eq:exp-nu-i-2}
\E[\nus]{f} = \frac{1}{3} \sum_{\sigma\in \{0,1,\star\}^V}\sum_{x\in\{0,1,\star\} } \nus(\sigma^x)f(\sigma^x) = \frac{1}{3} \sum_{\sigma\in \{0,1,\star\}^V}\sum_{x\in\{0,1,\star\} } \nu(\sigma^x)\cdot g(\sigma^x) =  \E[\nu]{g}.
\end{align}
In the last equation, we use the fact that $\nu$ is a distribution over $\Omega(\pi)$.
%Note that $f$ is an increasing function. We have the following claim about the function $g$.
%\begin{claim}\label{claim:g_i-increasing}
%The function $g: \Omega(\pi)\to \mathbb{R}$ is also increasing.
%\end{claim}
%We will prove \Cref{claim:g_i-increasing} later. Now, assuming \Cref{claim:g_i-increasing}, we continue to prove the lemma.
%Using the assumption that $\eta \preceq \nu$ and \Cref{claim:g_i-increasing}, we have
%\begin{align}\label{eq:ineq-1}
%\E[\tilde{\nu}^i]{f}=\E[\nu]{g_i}\geq \E[\eta]{g_i}.    
%\end{align}
Next, we claim the following relation between $\E[\nu]{g}$ and $\E[\nug]{f}$.
\begin{claim}\label{claim:g-i-f}
$\E[\nu]{g}\geq \E[\nug]{f}$.
\end{claim}
The proof of \Cref{claim:g-i-f} requires the condition that $\frac{\nu}{\pi}$ is an increasing function. The formal proof is given later.
Let us first assume that \Cref{claim:g-i-f} holds.
Combining the claim with~\eqref{eq:exp-nu-i-2},
\begin{align}\label{eq:ineq-i-each-2}
  \E[\nug]{f}\leq \E[\nus]{f}.
\end{align}
This verifies~\eqref{eq:ineq-i-each} and proves the lemma.
\end{proof}

\end{document}
